\documentclass[11pt]{article}

\usepackage[margin=1in]{geometry}
\usepackage{amsmath,amssymb,amsthm}
\usepackage{hyperref}
\usepackage{xurl}
\usepackage{microtype}
\usepackage{enumitem}

\hypersetup{
  colorlinks=true,
  linkcolor=blue,
  citecolor=blue,
  urlcolor=blue
}

\newtheorem{proposition}{Proposition}
\newtheorem{remark}{Remark}

\newcommand{\E}{\mathbb{E}}
\newcommand{\Var}{\mathrm{Var}}
\newcommand{\R}{\mathbb{R}}

\title{Preliminary Notes on ERM with Few Latent Tasks}
\author{}
\date{\today}

\begin{document}
\maketitle

\begin{abstract}
This note records preliminary findings for the following question. Data are generated hierarchically: first draw latent task variables
\[
\theta_1,\dots,\theta_k \overset{\mathrm{iid}}{\sim} p(\theta),
\]
then for each task draw
\[
Z_{ij} = (X_{ij},Y_{ij}) \overset{\mathrm{iid}}{\sim} P_{\theta_i},
\qquad j=1,\dots,m.
\]
We pool all $km$ examples and run empirical risk minimization (ERM) over a hypothesis class $\mathcal{H}$. When does the resulting predictor achieve low risk under the mixture distribution $D(z) = \int P_{\theta}(z)\,p(d\theta)$, especially in regimes where the number of distinct tasks $k$ is much smaller than the usual $O(\varepsilon^{-2})$ scale? The literature suggests that this problem sits between classical supervised learning, meta-learning, and domain generalization, but is not identical to any of them. The main preliminary conclusion is simple: pooled ERM is unbiased for the target mixture risk, but the real bottleneck is task-level heterogeneity. Large $m$ only helps with within-task noise; beating the $k$ bottleneck requires structural assumptions on the task-risk map $\theta \mapsto \E[\ell(h;Z)\mid \theta]$.
\end{abstract}

\section{Problem setup}
Let $\ell(h;z)$ be a loss in $[0,1]$. The population risk of a predictor $h \in \mathcal{H}$ is
\[
R(h) := \E_{\theta \sim p}\E_{Z \sim P_{\theta}}[\ell(h;Z)].
\]
Given hierarchical data $\{Z_{ij}\}_{i \in [k], j \in [m]}$, the pooled empirical risk is
\[
\widehat{R}_{k,m}(h) := \frac{1}{k}\sum_{i=1}^k \frac{1}{m}\sum_{j=1}^m \ell(h;Z_{ij}),
\]
and pooled ERM returns
\[
\widehat{h} \in \arg\min_{h \in \mathcal{H}} \widehat{R}_{k,m}(h).
\]

Two features distinguish this from ordinary iid learning.
\begin{enumerate}[label=(\roman*)]
\item The individual examples $Z_{ij}$ are not jointly iid because points from the same task share the same latent $\theta_i$.
\item The dependence is highly structured: the sample is a collection of $k$ iid clusters, each of size $m$.
\end{enumerate}

The question of interest is not whether $km$ examples are enough in the classical iid sense. Rather, it is whether a small number of distinct tasks can still suffice because the task family has special structure.

\section{Literature review}

\subsection{Motivation from in-context learning}
The motivating analogy is the line of work showing that in-context learning can recover learning algorithms that generalize beyond the finite task family seen during pretraining. Akyurek et al.~\cite{akyurek2023} showed that transformers can implement linear estimators such as gradient descent and ridge regression in context. Raventos et al.~\cite{raventos2023} sharpened the point by training transformers on a finite set of linear regression tasks and identifying a task-diversity threshold: below it, the transformer behaves like a Bayesian estimator tied to the pretraining task prior; above it, behavior aligns with ridge regression over a broader task family.

The present supervised-learning question is a natural analogue. Instead of asking whether a transformer trained on finitely many tasks can perform the right in-context algorithm on new tasks, we ask whether ordinary pooled ERM, trained on finitely many latent tasks, can recover a predictor with low risk under the full meta-distribution.

\subsection{Learning to learn from an environment of tasks}
Baxter's classical ``environment of tasks'' framework~\cite{baxter2000} is the closest foundational theory. The learner observes multiple tasks sampled from a task environment and uses them to learn an inductive bias that transfers to novel tasks. The central sample-complexity split is exactly the one relevant here: number of tasks versus number of examples per task.

Later work made this split more explicit. Maurer, Pontil, and Romera-Paredes~\cite{maurer2016} analyzed multitask and learning-to-learn representation learning, identifying regimes where many related tasks reduce the burden on per-task sample size. Denevi et al.~\cite{denevi2019} studied learning-to-learn with biased regularization and showed benefits when the variance across tasks is small. More recently, Aliakbarpour et al.~\cite{aliakbarpour2024} proved that in structured representation-learning settings, one can metalearn from very few samples per task, but only because the shared representation class is strongly restricted.

These papers do not analyze pooled ERM directly. Typically the algorithm learns a shared representation, bias vector, or initialization and then adapts to a fresh task. Still, they strongly suggest the right principle for the present problem: if one wants nontrivial generalization from few tasks, the family of task-level predictors or risks must be simple.

\subsection{Domain generalization and latent domains}
Another adjacent literature is domain generalization. Here multiple environments are available at training time and the goal is to perform well on unseen environments. IRM~\cite{arjovsky2019} and REx~\cite{krueger2021} search for predictors whose behavior is stable across environments. However, these methods usually assume observed environment labels and often target a stronger notion of robustness than the average-risk objective above.

For the current question, two facts from the domain-generalization literature matter. First, plain ERM is a strong baseline: Gulrajani and Lopez-Paz~\cite{gulrajani2020} showed that carefully tuned ERM is competitive with specialized domain-generalization methods on standard benchmarks, and Sener and Koltun~\cite{sener2023} argued that some domain-generalization penalties help only when they do not induce excess empirical risk. Second, there is work on latent domains, such as Matsuura and Harada~\cite{matsuura2019}, where domain labels are not observed and one must infer hidden environments.

That said, the objective here is still different. We do not seek worst-domain robustness, and we do not necessarily use domain labels at either train or test time. The target is the average mixture risk under $p(\theta)P_{\theta}$, with a clustered but exchangeable training sample.

\section{Preliminary findings}

\subsection{Unbiasedness is not the issue}
The first observation is almost trivial, but useful.

\begin{proposition}[Pooled empirical risk is unbiased]
For every fixed $h \in \mathcal{H}$,
\[
\E[\widehat{R}_{k,m}(h)] = R(h).
\]
\end{proposition}

\begin{proof}
Each individual point $Z_{ij}$ has marginal distribution
\[
\int P_{\theta}(z)\,p(d\theta) = D(z),
\]
so
\[
\E[\ell(h;Z_{ij})] = \E_{\theta}\E_{Z \sim P_{\theta}}[\ell(h;Z)] = R(h).
\]
Averaging over $i$ and $j$ preserves the same expectation.
\end{proof}

So the clustered sampling scheme does not introduce bias in the target objective. The problem is concentration: how quickly does $\widehat{R}_{k,m}(h)$ approximate $R(h)$, uniformly over $h$?

\subsection{Exact variance decomposition}
For a fixed $h$, define the task-conditional mean and variance
\[
\mu_h(\theta) := \E[\ell(h;Z)\mid \theta],
\qquad
v_h(\theta) := \Var(\ell(h;Z)\mid \theta).
\]

\begin{proposition}[Task-level versus within-task variance]
For every fixed $h \in \mathcal{H}$,
\[
\Var(\widehat{R}_{k,m}(h))
= \frac{1}{k}\Var_{\theta}(\mu_h(\theta))
  \;+\;
  \frac{1}{km}\E_{\theta}[v_h(\theta)].
\]
\end{proposition}

\begin{proof}
Let
\[
U_i(h) := \frac{1}{m}\sum_{j=1}^m \ell(h;Z_{ij}).
\]
The variables $U_1(h),\dots,U_k(h)$ are iid across $i$, and
\[
\widehat{R}_{k,m}(h) = \frac{1}{k}\sum_{i=1}^k U_i(h).
\]
Hence
\[
\Var(\widehat{R}_{k,m}(h)) = \frac{1}{k}\Var(U_1(h)).
\]
Applying the law of total variance and using conditional independence within a task,
\begin{align*}
\Var(U_1(h))
&= \Var\!\left(\E[U_1(h)\mid \theta_1]\right)
  + \E\!\left[\Var(U_1(h)\mid \theta_1)\right] \\
&= \Var_{\theta}(\mu_h(\theta))
  + \E_{\theta}\!\left[\frac{v_h(\theta)}{m}\right].
\end{align*}
Substituting gives the claim.
\end{proof}

This decomposition is the key preliminary finding. It says that $m$ only suppresses the within-task noise term. The task-level heterogeneity term, $\Var_{\theta}(\mu_h(\theta))/k$, is unaffected by increasing the number of samples from each already-observed task.

\paragraph{Design-effect form.}
Let
\[
\sigma_h^2 := \Var_{Z \sim D}(\ell(h;Z))
= \Var_{\theta}(\mu_h(\theta)) + \E_{\theta}[v_h(\theta)].
\]
Then
\[
\Var(\widehat{R}_{k,m}(h))
= \frac{\sigma_h^2}{km}\Bigl(1 + (m-1)\rho_h\Bigr),
\qquad
\rho_h := \frac{\Var_{\theta}(\mu_h(\theta))}{\sigma_h^2}.
\]
The quantity $\rho_h$ is the intraclass correlation of losses from the same task. Equivalently, the effective sample size is
\[
n_{\mathrm{eff}}(h) = \frac{km}{1 + (m-1)\rho_h}.
\]
When $\rho_h \approx 0$, pooled ERM behaves almost like iid learning with $km$ samples. When $\rho_h \approx 1$, only the number of tasks matters and the effective sample size is about $k$.

\subsection{A finite-class generalization bound}
The preceding variance formula converts immediately into a concentration bound for finite classes.

\begin{proposition}[Finite-class Bernstein bound]
Assume $\ell(h;z) \in [0,1]$ for all $h,z$, and let $\mathcal{H}$ be finite. Then with probability at least $1-\delta$, simultaneously for all $h \in \mathcal{H}$,
\[
\bigl|R(h) - \widehat{R}_{k,m}(h)\bigr|
\le
\sqrt{
\frac{2 \log(2|\mathcal{H}|/\delta)}{k}
\left(
\Var_{\theta}(\mu_h(\theta))
+ \frac{\E_{\theta}[v_h(\theta)]}{m}
\right)
}
+
\frac{\log(2|\mathcal{H}|/\delta)}{3k}.
\]
\end{proposition}

\begin{proof}
For fixed $h$, the task averages $U_i(h)$ from the previous proof are iid, bounded in $[0,1]$, with mean $R(h)$ and variance
\[
\Var(U_i(h)) = \Var_{\theta}(\mu_h(\theta)) + \frac{1}{m}\E_{\theta}[v_h(\theta)].
\]
Applying Bernstein's inequality to $\frac{1}{k}\sum_i U_i(h)$ and then taking a union bound over $\mathcal{H}$ yields the claim.
\end{proof}

This bound makes the bottleneck explicit. If the task-mean losses $\mu_h(\theta)$ vary substantially across $\theta$, then the first term dominates and one essentially needs many tasks. If instead the task-mean losses are nearly constant, then the bound is close to the ordinary iid $1/\sqrt{km}$ rate.

\subsection{Why small $k$ requires structure}
The previous propositions already imply an important impossibility heuristic.

\begin{remark}[The $m \to \infty$ reduction]
If $m$ were infinite, then for each observed task $\theta_i$ we would know the exact task risk $\mu_h(\theta_i)$ for every $h$. The learning problem would reduce to approximating
\[
R(h) = \E_{\theta}[\mu_h(\theta)]
\]
from the task sample $\theta_1,\dots,\theta_k$. Therefore, without additional structure on the function class
\[
\mathcal{M} := \{\mu_h : h \in \mathcal{H}\},
\]
large $m$ cannot compensate for small $k$.
\end{remark}

This is the cleanest way to state what seems both true and important:
\begin{enumerate}[label=(\alph*)]
\item A nontrivial ``few tasks suffice'' phenomenon is possible only if the task-level risk class $\mathcal{M}$ is simple, low-rank, low-variance, or otherwise compressible.
\item In a distribution-free setting over arbitrary task families, there is no miracle. Once within-task noise is averaged out, the remaining problem is ordinary generalization over the space of tasks.
\end{enumerate}

This is very much in line with the meta-learning literature: success with few samples per task is possible, but only because the shared structure across tasks is constrained.

\subsection{Linear regression as a concrete test case}
The linear-regression case is a natural place to start because it most closely parallels the ICL literature. Suppose
\[
h_w(x) = w^\top x,
\qquad
\ell(w;(x,y)) = (w^\top x - y)^2.
\]
For each task $\theta$, define
\[
A(\theta) := \E[XX^\top \mid \theta],
\qquad
b(\theta) := \E[XY \mid \theta],
\qquad
c(\theta) := \E[Y^2 \mid \theta].
\]
Then the task risk is
\[
r_{\theta}(w)
:= \E[(w^\top X - Y)^2 \mid \theta]
= w^\top A(\theta) w - 2 w^\top b(\theta) + c(\theta),
\]
and the mixture risk is
\[
R(w)
= w^\top \bar{A} w - 2 w^\top \bar{b} + \bar{c},
\]
where
\[
\bar{A} := \E_{\theta}[A(\theta)],
\qquad
\bar{b} := \E_{\theta}[b(\theta)],
\qquad
\bar{c} := \E_{\theta}[c(\theta)].
\]

This specialization exposes an appealing route for future work. The full meta-distribution matters only through the averages of a small set of task statistics. Therefore:
\begin{enumerate}[label=(\roman*)]
\item if $A(\theta)$ and $b(\theta)$ concentrate strongly across tasks, then small $k$ should suffice;
\item if the family $\{(A(\theta),b(\theta))\}$ lies near a low-dimensional subspace, then one should hope for sample complexity driven by that low task complexity rather than by an unrestricted $1/\varepsilon^2$ task count;
\item if $A(\theta)$ and $b(\theta)$ vary arbitrarily, then even $m = \infty$ does not avoid the task-sampling bottleneck.
\end{enumerate}

This is, in my view, the closest supervised analogue of the phenomenon isolated by Raventos et al.~\cite{raventos2023}: finitely many observed tasks can stand in for an infinite task family only when those tasks already pin down the relevant sufficient statistics of the broader environment.

\section{Takeaways and candidate directions}
At this stage, the picture seems to be the following.

\begin{enumerate}[label=(\arabic*)]
\item \textbf{Pooled ERM is not misaligned with the target risk.} The sample is dependent, but each example has the correct marginal distribution. The real issue is variance and uniform convergence, not bias.
\item \textbf{Task heterogeneity is the fundamental obstruction.} The exact variance decomposition shows that increasing $m$ only reduces within-task noise. It does nothing to reduce the variance of the task-mean losses across $\theta$.
\item \textbf{Any regime with very small $k$ must exploit structure on the task-risk map.} Promising structural assumptions include common invariant conditionals, low-rank families of task moments, shared sufficient statistics, or small task variance around a common predictor.
\item \textbf{Linear regression looks especially tractable.} One can hope for a theorem saying that if the task moments $(A(\theta),b(\theta))$ come from a low-complexity family, then pooled ERM or regularized ERM generalizes to the full mixture with $k$ much smaller than the naive iid task count.
\end{enumerate}

The main open question for the next step is therefore not merely ``can small $k$ work?'', but rather:
\[
\textit{what is the right notion of task-level complexity that replaces raw task count?}
\]
Finding the right answer here would give a supervised-learning counterpart to the finite-task-diversity phenomena already observed in in-context learning.

\begin{thebibliography}{11}

\bibitem{akyurek2023}
Ekin Akyurek, Dale Schuurmans, Jacob Andreas, Tengyu Ma, and Denny Zhou.
\newblock What learning algorithm is in-context learning? Investigations with linear models.
\newblock In \emph{International Conference on Learning Representations}, 2023.
\newblock \url{https://arxiv.org/abs/2211.15661}

\bibitem{raventos2023}
Allan Raventos, Mansheej Paul, Feng Chen, and Surya Ganguli.
\newblock Pretraining task diversity and the emergence of non-Bayesian in-context learning for regression.
\newblock arXiv:2306.15063, 2023.
\newblock \url{https://arxiv.org/abs/2306.15063}

\bibitem{baxter2000}
Jonathan Baxter.
\newblock A model of inductive bias learning.
\newblock \emph{Journal of Artificial Intelligence Research}, 12:149--198, 2000.
\newblock \url{https://arxiv.org/abs/1106.0245}

\bibitem{maurer2016}
Andreas Maurer, Massimiliano Pontil, and Bernardino Romera-Paredes.
\newblock The benefit of multitask representation learning.
\newblock \emph{Journal of Machine Learning Research}, 17(81):1--32, 2016.
\newblock \url{https://jmlr.org/beta/papers/v17/15-242.html}

\bibitem{denevi2019}
Giulia Denevi, Carlo Ciliberto, Riccardo Grazzi, and Massimiliano Pontil.
\newblock Learning-to-learn stochastic gradient descent with biased regularization.
\newblock arXiv:1903.10399, 2019.
\newblock \url{https://arxiv.org/abs/1903.10399}

\bibitem{aliakbarpour2024}
Maryam Aliakbarpour, Konstantina Bairaktari, Gavin Brown, Adam Smith, Nathan Srebro, and Jonathan Ullman.
\newblock Metalearning with very few samples per task.
\newblock arXiv:2312.13978, 2024.
\newblock \url{https://arxiv.org/abs/2312.13978}

\bibitem{arjovsky2019}
Martin Arjovsky, Leon Bottou, Ishaan Gulrajani, and David Lopez-Paz.
\newblock Invariant risk minimization.
\newblock arXiv:1907.02893, 2019.
\newblock \url{https://arxiv.org/abs/1907.02893}

\bibitem{krueger2021}
David Krueger, Ethan Caballero, J\"orn-Henrik Jacobsen, Amy Zhang, Jonathan Binas, Dinghuai Zhang, Bernhard Sch\"olkopf, and Aaron Courville.
\newblock Out-of-distribution generalization via risk extrapolation (REx).
\newblock In \emph{International Conference on Machine Learning}, 2021.
\newblock \url{https://arxiv.org/abs/2003.00688}

\bibitem{gulrajani2020}
Ishaan Gulrajani and David Lopez-Paz.
\newblock In search of lost domain generalization.
\newblock arXiv:2007.01434, 2020.
\newblock \url{https://arxiv.org/abs/2007.01434}

\bibitem{matsuura2019}
Toshihiko Matsuura and Tatsuya Harada.
\newblock Domain generalization using a mixture of multiple latent domains.
\newblock arXiv:1911.07661, 2019.
\newblock \url{https://arxiv.org/abs/1911.07661}

\bibitem{sener2023}
Ozan Sener and Vladlen Koltun.
\newblock Domain generalization without excess empirical risk.
\newblock arXiv:2308.15856, 2023.
\newblock \url{https://arxiv.org/abs/2308.15856}

\end{thebibliography}

\end{document}
